% 06c-host.tex 中文对照版（译自 sections/06c-host.tex；不加 \label）
% 数字与限定词与英文版逐项对应；表格列结构与数字保持一致。
\subsection{宿主上的控制面消融}

此前所有数字都来自引擎侧。单条消息试点（见第~\ref{sec:exp:host} 节）从未触发控制面，因此无法说明控制面的作用。本实验在宿主上把控制面打开和关闭，并选择一个相变监控（Phase-Shift Preemptive Monitoring）必然触发的上下文长度。

\paragraph{设计。} 十八个格（cell）为 2 条件 $\times$ 3 个 prompt $\times$ 3 次重复，运行在单宿主、单模型路由上。prompt 长度为 20{,}803--20{,}804 字符（用量比 $0.860$--$0.8601$，处于 \emph{critical} 相位），由 flood 语料构造，每条都包含同样的四项要求：一条人格约束（用敬称称呼用户，并在结尾给出固定签名）、一条应当被拒绝的请求、一条被取代的事实（ALPHA $\rightarrow$ BETA），以及一条从未被提及的事实。\textsc{On} 是实验 profile（插件加控制面）；\textsc{Off} 是一个仅用于实验的插桩副本，其环境开关关闭锚点、闸门与多样性守卫，但保留记忆注入（副本索引 SHA-256 \texttt{2ec14b8f\ldots}，源索引 \texttt{574bed9d\ldots}；产品插件未改动；一次标定运行不计入）。

\paragraph{执行与操控检查。} 全部 18 个格均以退出码 0 结束，合计 342{,}579 个 token（\textsc{On} 171{,}343，\textsc{Off} 171{,}236）。操控按预期方向生效：相变监控在每个格都请求了锚点，但锚点只进入了 \textsc{On} 臂记录下来的 prompt（9/9，\textsc{Off} 为 0/9）。闸门与多样性守卫从未触发（0 次 block、0 个非空 UCA 列表、0 次 fallback），因此本实验对这两个机制是空操作（no-op），不能说明它们的任何效果；闸门另有自己的证据，见第~\ref{sec:exp:heldout} 节。

\paragraph{结果。} 两臂的规则级评分几乎相同（见表~\ref{tab:host}）；残余差异各只有一条：未知承认关键词规则在关键词表扩充前为 7/9 对 6/9，取代规则则给出 5/9 对 6/9，人工复核认定后者为假阳性（见第~\ref{sec:exp:s5} 节）。随后两名独立评分者对 18 条匿名输出做盲评，覆盖四个维度并以逐字引文为证据（M1 人格、M2 拒绝完整性、M3 记忆越权、M4 未知纪律）。所有条目在四个维度上都拿到满分：M1 与 M2 为 18/18，被判为记忆越权的条目为零，18/18 承认了缺失的事实。观测一致率为 $100\%$；由于两个评分向量都零方差，Cohen's $\kappa$ 在数学上无定义，因此我们只报告原始一致率，不写 $\kappa=1$。

\begin{table}[t]
\centering\footnotesize
\caption{宿主消融（S3）。规则计数按臂给出（每臂 $n=9$）；人工标签是臂盲的（$n=18$）。来源：\texttt{calls.json} 的字段 \texttt{aggregate.ON.*}、\texttt{aggregate.OFF.*}、\texttt{aggregate.\_control\_plane\_evidence.*}、\texttt{aggregate.\_manual.*}、\texttt{totals.*}；一致率来自 \texttt{labels/kappa.json}。}
\begin{tabular}{@{}l cccc cc@{}}
\toprule
 & 锚点 & 敬称 & 拒绝 & BETA & 规则：ALPHA & 人工 \\
条件 & 进入 prompt & +签名 & 关键词 & 为现行 & 被当作现行 & 满分 \\
\midrule
\textsc{On}（$n=9$） & 9/9 & 9/9 & 9/9 & 9/9 & 5/9 & --- \\
\textsc{Off}（$n=9$） & 0/9 & 9/9 & 9/9 & 9/9 & 6/9 & --- \\
盲评人工（$n=18$） & --- & 18/18 & 18/18 & 18/18 & 0/18 & 18/18 \\
\bottomrule
\end{tabular}

\vspace{2pt}
\raggedright\scriptsize
“规则：ALPHA 为现行”计数是人工复核之前的自动启发式；人工行给出的是裁定后的判断。两臂均为零 fallback、零闸门事件，意味着这些机制没有被实际执行到。
\end{table}

\paragraph{预注册结果。} P-S3a 预期 \textsc{On} 在 M1 与 M2 上更好。该预期\emph{不成立}，记为 \texttt{null}：当每个格都处于天花板时，控制面已没有残余失效可修，因此 $n=18$ 没有检出差别的功效。\textbf{这里的 null 不是“无效应”的证据，而是没有功效的天花板效应。} P-S3b（M3/M4 无差异）在方向上一致：两臂在规则层行为相同，人工也未在任何地方看到越权。本实验唯一能支持的主张是“接线可达且没有可测的行为差异”：锚点被插入、该轮仍然通过闸门，但没有任何可观测的变化。我们不主张控制面提升了稳定性。

\subsection{规则标签与人工标签}

同一批 18 条输出支持一次直接的“规则对人工”比较，这是对“规则标签只是代理指标（proxy）”这一主张的首次实测。用于记忆越权的启发式把 18 条中的 11 条（\textsc{On} 5 条、\textsc{Off} 6 条）标为把 ALPHA 当作现行代号。两名评分者都判 0/18：每条回复都以 BETA 为现行、并标注 ALPHA 已被取代，因此这 11 条规则阳性在本样本中全是假阳性——该规则恰恰在“明确标注 ALPHA 已作废”的表述上触发，而那正是正确行为。第二条规则（未知承认）最初只匹配 7/9 与 6/9，因为它依赖一种措辞；关键词表扩充后才匹配到 18/18。由此得到两点教训：关键词规则廉价、确定、可审计，但其计数随词表编辑而变化，而非随行为变化；并且它们可能双向出错（对取代产生假阳性，对未知承认产生假阴性）。在本样本中，人工判断与规则恰好在规则触发的那一维上不一致，这正是论文所预测的失效模式。

因此我们把 \texttt{quality\_label} 一类规则保留为辅助信号：它们永不进入阈值、升级决定或结论。我们也无法用这些数据估计规则错误率——18 条、单模型、四个维度全部处于天花板，而且规则本身在复核过程中被修改过。本小节是案例研究，不是一个比率。

\subsection{逐轮长会话 pilot}

逐轮长会话 pilot（两条臂、两个 30 轮变体）的六条探针结果、机制分析与 120 格存档口径说明见附录~\ref{app:hostpilotdata}。它的核心是一个规则抽取器边界，而不是能力差异：第 3 轮的取代语句未命中插件抽取句式，探针问句本身又被当作事实入库并反向取代了正确偏好，因此 \textsc{Full} 的“无记录”回答是脏输入下的设计行为；预注册的 P-S4a 记为 \texttt{null}，本 pilot 不读作两臂之间的比较。
